\documentclass[tikz,border=4pt]{standalone}
\usepackage{ctex}
\usepackage{amsmath}
\usetikzlibrary{calc, arrows.meta}

\begin{document}
% thinking-toolkit/07 结构识别（换元法骨架图）
% 示意：ax^4+bx^2+c=0，令t=x^2，化为at^2+bt+c=0
\begin{tikzpicture}[scale=1.0, thick]
  % 原方程框
  \draw[rounded corners=4pt, fill=pink!30, draw=black]
    (-3.5, 0.5) rectangle (3.5, -0.5);
  \node[font=\small] at (0,0) {$ax^4 + bx^2 + c = 0$（四次方程）};

  % 换元箭头
  \draw[->,>=Stealth,thick] (0,-0.6) -- (0,-1.4)
    node[midway, right, font=\small] {令 $t = x^2\;(t\geq0)$};

  % 化简后框
  \draw[rounded corners=4pt, fill=orange!20, draw=black]
    (-3.0, -1.5) rectangle (3.0, -2.5);
  \node[font=\small] at (0,-2.0) {$at^2 + bt + c = 0$（关于$t$的二次方程）};

  % 解出t后还原
  \draw[->,>=Stealth,thick] (0,-2.6) -- (0,-3.4)
    node[midway, right, font=\small] {解出$t$，再由$x^2=t$求$x$};

  % 结果框
  \draw[rounded corners=4pt, fill=green!20, draw=black]
    (-3.0, -3.5) rectangle (3.0, -4.5);
  \node[font=\small] at (0,-4.0) {$x = \pm\sqrt{t}$（需检验 $t\geq0$）};

  \node[font=\footnotesize, gray] at (0,-5.1)
    {结构识别：隐藏的"二次结构"用换元法揭示};
\end{tikzpicture}
\end{document}
